\clearpage
\VerdierPhysicalPageAnchor{0190}%
\begingroup\small\setlength{\parskip}{0pt}\setlength{\emergencystretch}{1.5em}
\setlength{\abovedisplayskip}{4pt plus 2pt minus 2pt}%
\setlength{\belowdisplayskip}{4pt plus 2pt minus 2pt}%

\VerdierSemanticAnchor{ch3d-2-1-5}
\noindent\textbf{Remark 2.1.5.}\quad
Suppose that the hypotheses of Lemma (2.1.4) hold and, in addition, that
$X,Y\in\operatorname{Ob}\operatorname{comp}^{-}(\mathcal A)$. The argument
given above then shows that, after replacing $X$ if necessary by a
homotopy-equivalent complex, we may suppose that $m$ is an epimorphism.
The construction of the sequence $X_n$ given above then shows that the
inductive system $X_n$ is stationary (i.e. $X_n=Y$ and
$u_n=\operatorname{id}$ for small $n$) and that, for every $n$,
$X_n\in\operatorname{Ob}\operatorname{comp}^{-}(\mathcal A)$.

\VerdierSemanticAnchor{ch3d-2-2}
\subsection{Study of certain sets of objects.}

\VerdierSemanticAnchor{ch3d-2-2-1}
\noindent\textbf{2.2.1.}\quad
Let $M$ be a set of objects of $\mathcal A$. We shall consider the following
properties:

\VerdierSemanticAnchor{ch3d-2-2-1-e1}
\noindent$(E_1)$ Zero objects belong to $M$. The direct sum of two objects
belonging to $M$ is an object belonging to $M$.

\VerdierSemanticAnchor{ch3d-2-2-1-e2}
\noindent$(E_2)$ For every epimorphism $X\xrightarrow{u}L$, where $X$ is
an object of $\mathcal A$ and $L$ is an element of $M$, there is an
epimorphism $v\colon L'\longrightarrow L$, with $L'\in M$, that factors
through $u$.

\VerdierSemanticAnchor{ch3d-2-2-1-e3}
\noindent$(E_3)$ For every integer $n$ and every exact sequence
\[
0\longrightarrow X\longrightarrow L_n\longrightarrow L_{n-1}
\longrightarrow\cdots\longrightarrow L_0\longrightarrow0,
\qquad L_i\in M,\quad0\leq i\leq n,
\]
there is an epimorphism $L'\longrightarrow X$, where $L'\in M$.

\VerdierSemanticAnchor{ch3d-2-2-1-e4}
\noindent$(E_4)$ For every epimorphism
$u\colon L\longrightarrow L'$, with $L,L'\in M$, the kernel of $u$ is an
element of $M$.

\VerdierSemanticAnchor{ch3d-2-2-1-e5}
\noindent$(E_5)$ For every extension
$0\longrightarrow L\longrightarrow Z\longrightarrow L'\longrightarrow0$,
with $L,L'\in M$, the object $Z$ is an element of $M$; $0$ is an element
of $M$.

\VerdierSemanticAnchor{ch3d-2-2-1-e6}
\noindent$(E_6)$ The set $M$ is exact, i.e.\ whenever two of the objects in
a short exact sequence in $\mathcal A$ are elements of $M$, so is the
third; $0$ is an element of $M$.

\VerdierSemanticAnchor{ch3d-2-2-1-implications}
Note the implications $(E_6)\Rightarrow(E_4)$ and $(E_5)$,
$(E_4)\Rightarrow(E_3)$, and $(E_5)\Rightarrow(E_1)$. An unbounded exact
sequence

\VerdierSemanticAnchor{ch3d-2-2-1-infinite-presentation}
\[
\cdots\longrightarrow L_n\longrightarrow L_{n-1}
\longrightarrow\cdots\longrightarrow L_0\longrightarrow X
\longrightarrow0,
\qquad L_i\in M,\quad i\in\mathbb N,
\]
is called an $\infty$-$M$-presentation of $X$. Let $M$ be a set of objects
satisfying properties $(E_1)$, $(E_2)$, and $(E_3)$. Denote by
$M^{\wedge}$ the following set of objects: an object $X$ belongs to
$M^{\wedge}$ if it has an $\infty$-$M$-presentation and if, for every
integer $n$ and every exact sequence

\VerdierSemanticAnchor{ch3d-2-2-1-m-wedge}
\[
0\longrightarrow X'\longrightarrow L_n\longrightarrow L_{n-1}
\longrightarrow\cdots\longrightarrow L_0\longrightarrow X
\longrightarrow0,
\qquad L_i\in M,\quad0\leq i\leq n,
\]
the object $X'$ has an $\infty$-$M$-presentation.
\endgroup
